
tenu de la structure de @@M0@@ sur @@M1@@, déterminée par 4.7, et du fait que les @@M2@@ sont les classes des faisceaux intervenant dans le scindage de l'image inverse de @@M3@@ sur @@M4@@. Cela achève la démonstration de la première forme de 4.6.

Pour en déduire la deuxième forme, on remarque que les éléments @@M5@@ de la @@M6@@-@@M7@@-algèbre @@M8@@ vérifient les relations i) et ii). Si @@M9@@ est la @@M10@@-@@M11@@-algèbre engendrée par des éléments @@M12@@ satisfaisant i) et ii), il existe donc un @@M13@@-@@M14@@-homomorphisme de @@M15@@ dans @@M16@@ envoyant @@M17@@ sur @@M18@@ (V 4.7). D'après 4.6, cet homomorphisme est surjectif. Comme le @@M19@@-idéal engendré par les relations i) et ii) dans @@M20@@ contient l'idéal usuel engendré par les relations de 4.6, cet homomorphisme est aussi injectif, donc est un isomorphisme.

\textbf{Corollaire 4.8.} \uline{Avec les hypothèses et les notations de} 4.5, \uline{posons}~:
@@M21@@
\uline{alors la} @@M22@@-\uline{algèbre} @@M23@@ \uline{est engendrée par les} @@M24@@; \uline{l'idéal des relations entre les} @@M25@@ \uline{est engendré par les relations déduites de}~:
@@M26@@

Ce n'est qu'un autre énoncé de 4.6. On peut en effet exprimer les @@M27@@ en fonction des @@M28@@, et réciproquement. De plus, la relation~: @@M29@@ équivaut à la relation~: @@M30@@.

4.9.
Soit @@M31@@ un schéma; notons @@M32@@ la grassmannienne de type @@M33@@ de @@M34@@, avec @@M35@@. Pour @@M36@@ variable, les @@M37@@ forment un système inductif, le morphisme @@M38@@ étant défini en considérant @@M39@@ comme
